\chapter{2013年陕西卷（理）}

\section{单选题}

\begin{problem}
设全集为 \(\R\)，函数 \(f(x) = \sqrt{1 - x^2}\) 的定义域为 \(M\)，则 \(\complement_{\R} M\) 为（\quad）

\choices
    {\(\left[-1, 1\right]\)}
    {\(\left(-1, 1\right)\)}
    {\(\left(-\infty, -1\right] \cup \left[1, +\infty\right)\)}
    {\(\left(-\infty, -1\right) \cup \left(1, +\infty\right)\)}

\begin{answer}
D
\end{answer}

\begin{solution}
根式有意义的条件是 \(1-x^2\geqslant0\)，即 \(-1\leqslant x\leqslant1\)，所以 \(M=\left[-1,1\right]\). 因此 \(\complement_{\R}M=\left(-\infty,-1\right)\cup\left(1,+\infty\right)\)，故选 D.
\end{solution}
\end{problem}

\begin{problem}
根据下列算法语句，当输入 \(x\) 为 \(60\) 时，输出 \(y\) 的值为（\quad）

\bitmapfigure[max width=0.34\linewidth]{img/2013/shaanxi_science/q2_fig1.png}

\choices
    {\(25\)}
    {\(30\)}
    {\(31\)}
    {\(61\)}

\begin{answer}
C
\end{answer}

\begin{solution}
输入 \(x=60>50\)，程序执行 Else 分支，因此 \(y=25+0.6\times(60-50)=25+6=31\)，故选 C.
\end{solution}
\end{problem}

\begin{problem}
设 \(\bs{a}\)，\(\bs{b}\) 为向量，则 \(\left|\bs{a}\cdot\bs{b}\right|=\left|\bs{a}\right|\left|\bs{b}\right|\) 是“\(\bs{a}\parallel\bs{b}\)”的（\quad）

\choices
    {充分不必要条件}
    {必要不充分条件}
    {充分必要条件}
    {既不充分也不必要条件}

\begin{answer}
C
\end{answer}

\begin{solution}
若 \(\bs{a}\) 或 \(\bs{b}\) 是零向量，则两向量平行，且等式两边都为 \(0\). 若两向量都不是零向量，设它们的夹角为 \(\theta\)，则 \(\left|\bs{a}\cdot\bs{b}\right|=\left|\bs{a}\right|\left|\bs{b}\right|\left|\cos\theta\right|\)，所以题设等式成立当且仅当 \(\left|\cos\theta\right|=1\)，即 \(\theta=0\) 或 \(\pi\)，这也等价于 \(\bs{a}\parallel\bs{b}\). 因而二者互为充分必要条件，故选 C.
\end{solution}
\end{problem}

\begin{problem}
某单位有 \(840\) 名职工，现采用系统抽样方法，抽取 \(42\) 人做问卷调查，将 \(840\) 人按 \(1\)，\(2\)，\(\cdots\)，\(840\) 随机编号，则抽取的 \(42\) 人中，编号落入区间 \([481,720]\) 的人数为（\quad）

\choices
    {\(11\)}
    {\(12\)}
    {\(13\)}
    {\(14\)}

\begin{answer}
B
\end{answer}

\begin{solution}
系统抽样的间隔为 \(840\div42=20\). 设随机确定的第一个样本编号为 \(r\)，其中 \(1\leqslant r\leqslant20\)，则每个样本组抽取的编号为 \(r+20(k-1)\)，其中 \(k=1\)，\(2\)，\(\ldots\)，\(42\). 区间 \([481,720]\) 恰好由 \([481,500]\)，\([501,520]\)，\(\ldots\)，\([701,720]\) 这 \(12\) 个连续的长度为 \(20\) 的编号组组成，每组恰抽取 1 人，所以人数为 \(12\)，故选 B.
\end{solution}
\end{problem}

\begin{problem}
如图，在矩形区域 \(ABCD\) 的 \(A\)，\(C\) 两点处各有一个通信基站，假设其信号覆盖范围分别是扇形区域 \(ADE\) 和扇形区域 \(CBF\)（该矩形区域内无其他信号来源，基站工作正常）. 若在该矩形区域内随机地选一地点，则该地点无信号的概率是（\quad）

\bitmapfigure[max width=0.40\linewidth]{img/2013/shaanxi_science/q5_fig1.png}

\choices
    {\(1-\dfrac{\pi}{4}\)}
    {\(\dfrac{\pi}{2}-1\)}
    {\(2-\dfrac{\pi}{2}\)}
    {\(\dfrac{\pi}{4}\)}

\begin{answer}
A
\end{answer}

\begin{solution}
由图可知矩形的长、宽分别为 \(2\)，\(1\)，所以面积为 \(2\). 两个信号覆盖区域均为半径 \(1\) 的四分之一圆，且除边界外不重叠，覆盖总面积为 \(2\times\dfrac{\pi}{4}=\dfrac{\pi}{2}\). 因此无信号区域的面积为 \(2-\dfrac{\pi}{2}\)，所求概率为 \(\dfrac{2-\frac{\pi}{2}}{2}=1-\dfrac{\pi}{4}\)，故选 A.
\end{solution}
\end{problem}

\begin{problem}
设 \(z_{1}\)，\(z_{2}\) 是复数，则下列命题中的假命题是（\quad）

\choices
    {若 \(\left|z_1-z_2\right|=0\)，则 \(\overline{z_1}=\overline{z_2}\)}
    {若 \(z_1=\overline{z_2}\)，则 \(\overline{z_1}=z_2\)}
    {若 \(\left|z_1\right|=\left|z_2\right|\)，则 \(z_1\cdot\overline{z_1}=z_2\cdot\overline{z_2}\)}
    {若 \(\left|z_1\right|=\left|z_2\right|\)，则 \(z_1^2=z_2^2\)}

\begin{answer}
D
\end{answer}

\begin{solution}
选项 A 中，\(\left|z_1-z_2\right|=0\) 得 \(z_1=z_2\)，共轭后仍有 \(\overline{z_1}=\overline{z_2}\). 选项 B 两边取共轭即可得到结论；选项 C 中 \(z_i\overline{z_i}=\left|z_i\right|^2\)，所以结论成立. 对选项 D，取 \(z_1=1\)，\(z_2=\i\)，则 \(\left|z_1\right|=\left|z_2\right|=1\)，但 \(z_1^2=1\)，\(z_2^2=-1\)，结论不成立，故选 D.
\end{solution}
\end{problem}

\begin{problem}
设 \(\triangle ABC\) 的内角 \(A\)，\(B\)，\(C\) 所对的边分别为 \(a\)，\(b\)，\(c\)，若 \(b\cos C+c\cos B=a\sin A\)，则 \(\triangle ABC\) 的形状为（\quad）

\choices
    {锐角三角形}
    {直角三角形}
    {钝角三角形}
    {不确定}

\begin{answer}
B
\end{answer}

\begin{solution}
由余弦定理，\(b\cos C=\dfrac{a^2+b^2-c^2}{2a}\)，\(c\cos B=\dfrac{a^2+c^2-b^2}{2a}\). 两式相加得 \(b\cos C+c\cos B=a\). 代入题设并注意 \(a>0\)，得 \(a=a\sin A\)，即 \(\sin A=1\). 因为 \(0<A<\pi\)，所以 \(A=\dfrac{\pi}{2}\)，三角形为直角三角形，故选 B.
\end{solution}
\end{problem}

\begin{problem}
设函数 \(f(x)=\begin{cases}
\left(x-\dfrac{1}{x}\right)^6, & x<0,\\
-\sqrt{x}, & x\geqslant0,
\end{cases}\) 则当 \(x>0\) 时，\(f[f(x)]\) 表达式的展开式中常数项为（\quad）

\choices
    {\(-20\)}
    {\(20\)}
    {\(-15\)}
    {\(15\)}

\begin{answer}
A
\end{answer}

\begin{solution}
当 \(x>0\) 时，\(f(x)=-\sqrt{x}<0\)，因此应使用函数的第一段：\(f[f(x)]=\left(-\sqrt{x}-\dfrac{1}{-\sqrt{x}}\right)^6=\left(\dfrac{1}{\sqrt{x}}-\sqrt{x}\right)^6=x^{-3}(1-x)^6\). 展开得 \(x^{-3}\sum_{k=0}^{6}(-1)^k\mathrm C_6^k x^k\). 常数项对应 \(k=3\)，其系数为 \((-1)^3\mathrm C_6^3=-20\)，故选 A.
\end{solution}
\end{problem}

\begin{problem}
在如图所示的锐角三角形空地中，欲建一个面积不小于 \(300 \mathrm{~m}^{2}\) 的内接矩形花园（阴影部分），则其边长 \(x\)（单位 \(\mathrm{m}\)）的取值范围是（\quad）

\bitmapfigure[max width=0.36\linewidth]{img/2013/shaanxi_science/q9_fig1.png}

\choices
    {\(\left[15,20\right]\)}
    {\(\left[12,25\right]\)}
    {\(\left[10,30\right]\)}
    {\(\left[20,30\right]\)}

\begin{answer}
C
\end{answer}

\begin{solution}
设矩形的另一边长为 \(y\). 三角形的底边和高均为 \(40\)，由相似三角形可得矩形上边所在的水平截面长为 \(40-y\)，所以 \(x=40-y\)，即 \(y=40-x\). 矩形面积不小于 \(300\) 给出 \(x(40-x)\geqslant300\)，整理得 \(x^2-40x+300\leqslant0\)，即 \((x-10)(x-30)\leqslant0\). 因而 \(10\leqslant x\leqslant30\)，故选 C.
\end{solution}
\end{problem}

\begin{problem}
设 \([x]\) 表示不大于 \(x\) 的最大整数，则对任意实数 \(x\)，\(y\)，有（\quad）

\choices
    {\(\left[-x\right]=-\left[x\right]\)}
    {\(\left[2x\right]=2\left[x\right]\)}
    {\(\left[x+y\right]\leqslant\left[x\right]+\left[y\right]\)}
    {\(\left[x-y\right]\leqslant\left[x\right]-\left[y\right]\)}

\begin{answer}
D
\end{answer}

\begin{solution}
选项 A 取 \(x=\dfrac12\)，有 \(\left[-x\right]=-1\neq0=-\left[x\right]\). 选项 B 取 \(x=\dfrac12\)，有 \(\left[2x\right]=1\neq0=2\left[x\right]\). 选项 C 取 \(x=y=\dfrac12\)，有 \(\left[x+y\right]=1>0=\left[x\right]+\left[y\right]\). 对选项 D，设 \(\left[x\right]=m\)，\(\left[y\right]=n\)，则可写成 \(x=m+\alpha\)，\(y=n+\beta\)，其中 \(0\leqslant\alpha\)，\(\beta<1\). 于是 \(\left[x-y\right]=m-n+\left[\alpha-\beta\right]\)，而 \(-1<\alpha-\beta<1\)，故 \(\left[\alpha-\beta\right]\leqslant0\)，从而 \(\left[x-y\right]\leqslant m-n=\left[x\right]-\left[y\right]\)，故选 D.
\end{solution}
\end{problem}

\section{填空题}

\begin{problem}
双曲线 \(\dfrac{x^2}{16}-\dfrac{y^2}{m}=1\) 的离心率为 \(\dfrac54\)，则 \(m\) 等于\fillinblank{}.

\begin{answer}
\(9\)
\end{answer}

\begin{solution}
双曲线的半实轴长满足 \(a^2=16\)，所以 \(a=4\). 由离心率公式 \(e=\dfrac ca=\dfrac54\) 得 \(c=5\). 又 \(c^2=a^2+b^2=16+m\)，所以 \(25=16+m\)，即 \(m=9\)，故填 \(9\).
\end{solution}
\end{problem}

\begin{problem}
某几何体的三视图如图所示，则其体积为\fillinblank{}.

\bitmapfigure[max width=0.34\linewidth]{img/2013/shaanxi_science/q12_fig1.png}

\begin{answer}
\(\dfrac{\pi}{3}\)
\end{answer}

\begin{solution}
由三视图可知，该几何体是底面半径为 \(1\) 的半个圆锥，圆锥的高为 \(2\). 因而其体积为 \(\dfrac12\times\dfrac13\pi\times1^2\times2=\dfrac\pi3\)，故填 \(\dfrac\pi3\).
\end{solution}
\end{problem}

\begin{problem}
若点 \((x,y)\) 位于曲线 \(y=\left|x-1\right|\) 与 \(y=2\) 所围成的封闭区域，则 \(2x-y\) 的最小值为\fillinblank{}.

\begin{answer}
\(-4\)
\end{answer}

\begin{solution}
封闭区域是由 \(y=\left|x-1\right|\) 和 \(y=2\) 围成的三角形，其三个顶点为 \((-1,2)\)，\((1,0)\)，\((3,2)\). 线性函数 \(2x-y\) 在该三角形上的最小值出现在顶点处，分别为 \(-4\)，\(2\)，\(4\)，所以最小值为 \(-4\)，故填 \(-4\).
\end{solution}
\end{problem}

\begin{problem}
观察下列等式：\(1^2=1\)；\(1^2-2^2=-3\)；\(1^2-2^2+3^2=6\)；\(1^2-2^2+3^2-4^2=-10\). 照此规律，第 \(n\) 个等式可为\fillinblank{}.

\begin{answer}
\(1^2-2^2+3^2-\cdots+(-1)^{n-1}n^2=(-1)^{n-1}\dfrac{n(n+1)}{2}\)
\end{answer}

\begin{solution}
记左端前 \(n\) 项的和为 \(S_n=\sum_{k=1}^{n}(-1)^{k-1}k^2\). 当 \(n=1\) 时，\(S_1=1=(-1)^0\dfrac{1\cdot2}{2}\). 假设 \(S_n=(-1)^{n-1}\dfrac{n(n+1)}{2}\)，则 \(S_{n+1}=S_n+(-1)^n(n+1)^2=(-1)^n\left[-\dfrac{n(n+1)}2+(n+1)^2\right]=(-1)^n\dfrac{(n+1)(n+2)}2\). 由数学归纳法，\(S_n=(-1)^{n-1}\dfrac{n(n+1)}2\)，故所填等式为 \(1^2-2^2+3^2-\cdots+(-1)^{n-1}n^2=(-1)^{n-1}\dfrac{n(n+1)}2\).
\end{solution}
\end{problem}

\section{选做题}

\begin{problem}
已知 \(a\)，\(b\)，\(m\)，\(n\) 均为正数，且 \(a+b=1\)，\(mn=2\)，则 \((am+bn)(bm+an)\) 的最小值为\fillinblank{}.

\begin{answer}
\(2\)
\end{answer}

\begin{solution}
展开并配方得 \((am+bn)(bm+an)=ab(m-n)^2+mn(a+b)^2=ab(m-n)^2+2\geqslant2\). 当 \(m=n=\sqrt2\) 时满足 \(mn=2\)，且等号成立，所以最小值为 \(2\)，故填 \(2\).
\end{solution}
\end{problem}

\begin{problem}
如图，弦 \(AB\) 与 \(CD\) 相交于 \(\odot O\) 内一点 \(E\)，过 \(E\) 作 \(BC\) 的平行线与 \(AD\) 的延长线相交于点 \(P\). 已知 \(PD=2DA=2\)，则 \(PE=\fillinblank{}\).

\bitmapfigure[max width=0.34\linewidth]{img/2013/shaanxi_science/q16_fig1.png}

\begin{answer}
\(\sqrt6\)
\end{answer}

\begin{solution}
因为 \(P\)，\(D\)，\(A\) 共线且 \(D\)，\(A\) 在 \(P\) 的同一侧，所以 \(\angle EPD=\angle EPA\). 又因 \(PE\parallel BC\)，\(ED\) 与 \(CD\) 共线，\(EA\) 与 \(AB\) 共线，有 \(\angle PED=\angle BCD\)，\(\angle PAE=\angle DAB\). 点 \(A\)，\(B\)，\(C\)，\(D\) 共圆，故 \(\angle BCD=\angle DAB\)，从而 \(\angle PED=\angle PAE\). 因此 \(\triangle PED\sim\triangle PEA\)，对应边满足 \(\dfrac{PE}{PA}=\dfrac{PD}{PE}\)，即 \(PE^2=PA\cdot PD\). 由 \(PD=2DA=2\) 得 \(DA=1\)，从而 \(PA=PD+DA=3\). 所以 \(PE^2=3\times2=6\)，又 \(PE>0\)，故 \(PE=\sqrt6\).
\end{solution}
\end{problem}

\begin{problem}
如图，以过原点的直线的倾斜角 \(\theta\) 为参数，则圆 \(x^2+y^2-x=0\) 的参数方程为\fillinblank{}.

\bitmapfigure[max width=0.30\linewidth]{img/2013/shaanxi_science/q17_fig1.png}

\begin{answer}
\(x=\cos^2\theta\)，\(y=\sin\theta\cos\theta\)
\end{answer}

\begin{solution}
过原点的直线斜率为 \(\tan\theta\)，故非竖直时可写为 \(y=x\tan\theta\). 代入圆方程得 \(x^2(1+\tan^2\theta)-x=0\)，除去原点对应的交点后，另一个交点满足 \(x=\dfrac{1}{1+\tan^2\theta}=\cos^2\theta\). 于是 \(y=x\tan\theta=\cos^2\theta\tan\theta=\sin\theta\cos\theta\). 当 \(\cos\theta=0\) 时，上式给出原点，也满足圆方程，因此参数方程为 \(x=\cos^2\theta\)，\(y=\sin\theta\cos\theta\).
\end{solution}
\end{problem}

\section{解答题}

\begin{problem}
已知向量 \(\bs{a}=\left(\cos x,-\dfrac12\right)\)，\(\bs{b}=\left(\sqrt3\sin x,\cos2x\right)\)，\(x\in\R\)，设函数 \(f(x)=\bs{a}\cdot\bs{b}\).

\begin{enumerate}
\item 求 \(f(x)\) 的最小正周期.
\item 求 \(f(x)\) 在 \(\left[0,\dfrac\pi2\right]\) 上的最大值和最小值.
\end{enumerate}

\end{problem}

\begin{solution}
\begin{enumerate}
\item \(f(x)=\sqrt3\sin x\cos x-\dfrac12\cos2x=\dfrac{\sqrt3}{2}\sin2x-\dfrac12\cos2x=\sin\left(2x-\dfrac\pi6\right)\). 因此 \(f(x)\) 的最小正周期为 \(\pi\).
\item 当 \(x\in\left[0,\dfrac\pi2\right]\) 时，\(2x-\dfrac\pi6\in\left[-\dfrac\pi6,\dfrac{5\pi}6\right]\). 在这个区间内，正弦函数的最大值为 \(1\)，且在 \(2x-\dfrac\pi6=\dfrac\pi2\)，即 \(x=\dfrac\pi3\) 时取得. 区间左端点给出 \(f(0)=\sin\left(-\dfrac\pi6\right)=-\dfrac12\)，右端点给出 \(f\left(\dfrac\pi2\right)=\sin\dfrac{5\pi}6=\dfrac12\)，所以最小值为 \(-\dfrac12\)，最大值为 \(1\).
\end{enumerate}
\end{solution}

\begin{problem}
设 \(\{a_n\}\) 是公比为 \(q\) 的等比数列.

\begin{enumerate}
\item 推导 \(\{a_n\}\) 的前 \(n\) 项和公式；
\item 设 \(q\neq1\)，证明数列 \(\{a_n+1\}\) 不是等比数列.
\end{enumerate}

\end{problem}

\begin{solution}
\begin{enumerate}
\item 设前 \(n\) 项和为 \(S_n\). 当 \(q=1\) 时，\(S_n=na_1\). 当 \(q\neq1\) 时，\(S_n=a_1+a_1q+\cdots+a_1q^{n-1}\)，两边乘以 \(q\) 得 \(qS_n=a_1q+a_1q^2+\cdots+a_1q^n\). 两式相减得 \((1-q)S_n=a_1(1-q^n)\)，所以 \(S_n=\dfrac{a_1(1-q^n)}{1-q}=\dfrac{a_1(q^n-1)}{q-1}\).
\item 反设 \(\{a_n+1\}\) 是等比数列. 取其前三项，则有 \((a_2+1)^2=(a_1+1)(a_3+1)\). 由于 \(a_2=a_1q\)，\(a_3=a_1q^2\)，代入并化简得 \(a_1(q-1)^2=0\). 等比数列的首项 \(a_1\neq0\)，所以 \(q=1\)，与已知 \(q\neq1\) 矛盾. 因此 \(\{a_n+1\}\) 不是等比数列.
\end{enumerate}
\end{solution}

\begin{problem}
如图，四棱柱 \(ABCD-A_{1}B_{1}C_{1}D_{1}\) 的底面 \(ABCD\) 是正方形，\(O\) 为底面中心，\(A_{1}O\perp\) 平面 \(ABCD\)，\(AB=AA_{1}=\sqrt2\).

\begin{enumerate}
\item 证明：\(A_{1}C\perp\) 平面 \(BB_{1}D_{1}D\).
\item 求平面 \(OCB_{1}\) 与平面 \(BB_{1}D_{1}D\) 的夹角 \(\theta\) 的大小.
\end{enumerate}

\bitmapfigure[max width=0.62\linewidth]{img/2013/shaanxi_science/q20_fig1.png}

\end{problem}

\begin{solution}
建立空间直角坐标系，使 \(O\) 为原点，\(OA\)，\(OB\)，\(OA_1\) 分别为三个坐标轴的正向. 因为正方形边长为 \(\sqrt2\)，所以其对角线长为 \(2\)，从而 \(OA=OB=1\). 又在直角三角形 \(AOA_1\) 中，\(AA_1=\sqrt2\)，故 \(OA_1=1\). 因此可取
\(A=(1,0,0)\)，\(B=(0,1,0)\)，\(C=(-1,0,0)\)，\(D=(0,-1,0)\)，\(A_1=(0,0,1)\). 由四棱柱的平移关系，\(B_1=(-1,1,1)\)，\(D_1=(-1,-1,1)\).

\begin{enumerate}
\item \(\overrightarrow{A_1C}=(-1,0,-1)\)，\(\overrightarrow{BB_1}=(-1,0,1)\)，\(\overrightarrow{BD}=(0,-2,0)\). 因此 \(\overrightarrow{A_1C}\cdot\overrightarrow{BB_1}=0\)，且 \(\overrightarrow{A_1C}\cdot\overrightarrow{BD}=0\). 直线 \(BB_1\) 与 \(BD\) 是平面 \(BB_1D_1D\) 内的两条相交直线，所以 \(A_1C\perp\) 平面 \(BB_1D_1D\).
\item 平面 \(OCB_1\) 的两个方向向量可取 \(\overrightarrow{OC}=(-1,0,0)\) 和 \(\overrightarrow{OB_1}=(-1,1,1)\)，故其法向量可取 \(\bs n_1=(0,1,-1)\). 平面 \(BB_1D_1D\) 的两个方向向量可取 \(\overrightarrow{BB_1}=(-1,0,1)\) 和 \(\overrightarrow{BD}=(0,-2,0)\)，故其法向量可取 \(\bs n_2=(1,0,1)\). 于是 \(\left|\bs n_1\cdot\bs n_2\right|=1\)，\(\left|\bs n_1\right|=\left|\bs n_2\right|=\sqrt2\)，所以 \(\cos\theta=\dfrac12\). 取两平面的锐角夹角，得 \(\theta=\dfrac\pi3\).
\end{enumerate}
\end{solution}

\begin{problem}
在一场娱乐晚会上，有 \(5\) 位民间歌手（\(1\) 至 \(5\) 号）登台演唱. 由现场数百名观众投票选出最受欢迎歌手. 各位观众须彼此独立地在选票上选 \(3\) 名歌手，其中观众甲是 \(1\) 号歌手的歌迷，他必选 \(1\) 号，不选 \(2\) 号，另在 \(3\) 至 \(5\) 号中随机选 \(2\) 名. 观众乙和丙对 \(5\) 位歌手的演唱没有偏爱，因此在 \(1\) 至 \(5\) 号中随机选 \(3\) 名歌手.

\begin{enumerate}
\item 求观众甲选中 \(3\) 号歌手且观众乙未选中 \(3\) 号歌手的概率；
\item \(X\) 表示 \(3\) 号歌手得到观众甲、乙、丙的票数之和，求 \(X\) 的分布列和数学期望.
\end{enumerate}

\end{problem}

\begin{solution}
\begin{enumerate}
\item 观众甲从 \(3\)，\(4\)，\(5\) 号中选取 \(2\) 名，等可能的选法有 \(\mathrm C_3^2=3\) 种，其中选中 \(3\) 号的选法有 \(2\) 种，所以观众甲选中 \(3\) 号的概率为 \(\dfrac23\). 观众乙从 \(5\) 位歌手中选取 \(3\) 名，选中 \(3\) 号的概率为 \(\dfrac{\mathrm C_4^2}{\mathrm C_5^3}=\dfrac35\)，所以未选中的概率为 \(\dfrac25\). 由于甲、乙彼此独立，所求概率为 \(\dfrac23\times\dfrac25=\dfrac4{15}\).
\item 设事件 \(A\)，\(B\)，\(C\) 分别表示观众甲、乙、丙选中 \(3\) 号歌手. 则 \(P(A)=\dfrac23\)，\(P(B)=P(C)=\dfrac35\)，且三个事件相互独立. 随机变量 \(X\) 的可能取值为 \(0\)，\(1\)，\(2\)，\(3\). 令 \(t\) 为形式变量，用 \(t^r\) 记录恰有 \(r\) 张票投给 \(3\) 号歌手. 于是三个观众对应的概率生成因子分别为：观众甲为 \(\dfrac13+\dfrac23t\)，观众乙和丙各为 \(\dfrac25+\dfrac35t\). 因而

\[
\left(\dfrac13+\dfrac23t\right)\left(\dfrac25+\dfrac35t\right)^2
=\dfrac4{75}+\dfrac4{15}t+\dfrac{11}{25}t^2+\dfrac6{25}t^3
\]

比较 \(t^r\) 的系数，得到 \(P(X=0)=\dfrac4{75}\)，\(P(X=1)=\dfrac4{15}\)，\(P(X=2)=\dfrac{11}{25}\)，\(P(X=3)=\dfrac6{25}\).

因此分布列为

\begin{center}
\begin{tabular}{c|cccc}
\(X\) & \(0\) & \(1\) & \(2\) & \(3\) \\
\hline
\(P\) & \(\dfrac4{75}\) & \(\dfrac4{15}\) & \(\dfrac{11}{25}\) & \(\dfrac6{25}\)
\end{tabular}
\end{center}

数学期望为 \(E(X)=0\cdot\dfrac4{75}+1\cdot\dfrac4{15}+2\cdot\dfrac{11}{25}+3\cdot\dfrac6{25}=\dfrac{28}{15}\).
\end{enumerate}
\end{solution}

\begin{problem}
已知动圆过定点 \(A(4,0)\)，且在 \(y\) 轴上截得的弦 \(MN\) 的长为 \(8\).

\begin{enumerate}
\item 求动圆圆心的轨迹 \(C\) 的方程；
\item 已知点 \(B(-1,0)\)，设不垂直于 \(x\) 轴的直线 \(l\) 与轨迹 \(C\) 交于不同的两点 \(P\)，\(Q\)，若 \(x\) 轴是 \(\angle PBQ\) 的角平分线，证明直线 \(l\) 过定点.
\end{enumerate}
\end{problem}

\begin{solution}
\begin{enumerate}
\item 设动圆圆心为 \(O_1(u,v)\)，半径为 \(R\). 设弦 \(MN\) 的中点为 \(H\)，则 \(O_1H\perp MN\)，且 \(H\) 在 \(y\) 轴上. 因为弦长为 \(8\)，所以 \(MH=4\)，从而 \(R^2=O_1H^2+MH^2=u^2+16\). 又圆过点 \(A(4,0)\)，故 \(R^2=(u-4)^2+v^2\). 两式相等得 \((u-4)^2+v^2=u^2+16\)，即 \(v^2=8u\). 反之，若 \(v^2=8u\)，以 \(O_1(u,v)\) 为圆心、\(R=\sqrt{u^2+16}\) 为半径作圆，则 \((u-4)^2+v^2=u^2+16=R^2\)，该圆过点 \(A(4,0)\)，且它在 \(y\) 轴上的弦长为 \(2\sqrt{R^2-u^2}=8\). 所以轨迹 \(C\) 的方程为 \(y^2=8x\).
\item 设直线 \(l\) 的方程为 \(y=kx+b\). 若 \(k=0\)，它与 \(y^2=8x\) 至多有一个交点，不符合题设，故 \(k\neq0\). 设 \(P(x_1,y_1)\)，\(Q(x_2,y_2)\)，则 \(x_1\)，\(x_2\) 是方程 \(k^2x^2+(2bk-8)x+b^2=0\) 的两个不同实根，所以 \(x_1+x_2=\dfrac{8-2bk}{k^2}\)，\(x_1x_2=\dfrac{b^2}{k^2}\).

\(x\) 轴是 \(\angle PBQ\) 的角平分线，关于 \(x\) 轴对称的两条射线斜率互为相反数，因此 \(\dfrac{y_1}{x_1+1}=-\dfrac{y_2}{x_2+1}\). 代入 \(y_i=kx_i+b\) 得 \(2kx_1x_2+(k+b)(x_1+x_2)+2b=0\). 再代入韦达定理，得到
\(2kb^2+(k+b)(8-2bk)+2bk^2=0\)，化简为 \(8(k+b)=0\)，故 \(b=-k\). 因而 \(l\) 的方程为 \(y=k(x-1)\)，恒过定点 \((1,0)\).
\end{enumerate}
\end{solution}

\begin{problem}
已知函数 \(f(x)=\e^x\)，\(x\in\R\).

\begin{enumerate}
\item 若直线 \(y=kx+1\) 与 \(f(x)\) 的反函数的图象相切，求实数 \(k\) 的值；
\item 设 \(x>0\)，讨论曲线 \(y=f(x)\) 与曲线 \(y=mx^2\)（\(m>0\)）公共点的个数；
\item 设 \(a<b\)，比较 \(\dfrac{f(a)+f(b)}{2}\) 与 \(\dfrac{f(b)-f(a)}{b-a}\) 的大小，并说明理由.
\end{enumerate}
\end{problem}

\begin{solution}
\begin{enumerate}
\item \(f(x)=\e^x\) 的反函数为 \(g(x)=\ln x\)，其中 \(x>0\). 设直线与曲线 \(y=\ln x\) 在 \(x=t\) 处相切，则切点满足 \(\ln t=kt+1\)，且切线斜率满足 \(k=g'(t)=\dfrac1t\). 因而 \(\ln t=2\)，得到 \(t=\e^2\)，所以 \(k=\e^{-2}\).
\item 由 \(\e^x=mx^2\) 得 \(m=\dfrac{\e^x}{x^2}\). 令 \(\varphi(x)=\dfrac{\e^x}{x^2}\)，则 \(\varphi'(x)=\dfrac{\e^x(x-2)}{x^3}\). 另外，\(\displaystyle\lim_{x\to0^+}\varphi(x)=+\infty\)，\(\displaystyle\lim_{x\to+\infty}\varphi(x)=+\infty\). 所以 \(\varphi\) 在 \((0,2)\) 上递减，在 \((2,+\infty)\) 上递增，且最小值为 \(\varphi(2)=\dfrac{\e^2}{4}\). 因此，方程 \(\varphi(x)=m\) 在 \(x>0\) 上：当 \(0<m<\dfrac{\e^2}{4}\) 时无解，即两曲线没有公共点；当 \(m=\dfrac{\e^2}{4}\) 时只有解 \(x=2\)，即有一个公共点；当 \(m>\dfrac{\e^2}{4}\) 时在 \((0,2)\) 和 \((2,+\infty)\) 各有一个解，即有两个公共点.
\item 函数 \(f(x)=\e^x\) 满足 \(f''(x)=\e^x>0\)，所以它在 \([a,b]\) 上严格凸. 对任意 \(x\in(a,b)\)，严格凸性给出

\[
f(x)<\dfrac{b-x}{b-a}f(a)+\dfrac{x-a}{b-a}f(b)
\]

对 \(x\) 从 \(a\) 到 \(b\) 积分，得到

\[
\int_a^b f(x)\,\mathrm dx
<\int_a^b\left(\dfrac{b-x}{b-a}f(a)+\dfrac{x-a}{b-a}f(b)\right)\mathrm dx
=\dfrac{b-a}{2}\bigl(f(a)+f(b)\bigr)
\]

由于 \(b-a>0\)，再除以 \(b-a\)，并利用 \(\int_a^b \e^x\,\mathrm dx=f(b)-f(a)\)，可得

\[
\dfrac{f(b)-f(a)}{b-a}<\dfrac{f(a)+f(b)}2
\]
\end{enumerate}
\end{solution}
